\documentclass[10pt,a4paper]{amsart}
\usepackage[margin=19mm]{geometry}
\usepackage{amsmath,amssymb,mathtools}
\usepackage[scr=rsfs,cal=euler]{mathalfa}
\usepackage{xfrac}
\usepackage[unicode,hidelinks,bookmarks=false]{hyperref}
\newcommand{\ie}{i.e.\ }
\newcommand{\cf}{cf.\ }
\newcommand{\resp}{respectively\ }
\newcommand{\Subsection}{\subsection}
\providecommand{\nobreakdash}{\nobreak\hbox{-}}
\setlength{\emergencystretch}{2em}
\multlinegap0pt
\usepackage{amssymb}
\usepackage{xcolor}
\usepackage[shortlabels]{enumitem}
\usepackage{mathtools}
\newtheorem{lemma}{Lemma}
\title{A discrepancy result for Hilbert modular forms}
\author{Baskar Balasubramanyam, Jishu Das, Kaneenika Sinha}
\date{Source: arXiv:2307.16736v3 (CC BY 4.0)}
\begin{document}
\maketitle

\noindent\emph{Source and licence.} A discrepancy result for Hilbert modular forms. Version arXiv:2307.16736v3 (CC BY 4.0), distributed by arXiv under the Creative Commons Attribution 4.0 International licence, http://creativecommons.org/licenses/by/4.0/, as recorded in the arXiv licence metadata for this exact version and shown on the abstract page https://arxiv.org/abs/2307.16736v3. Full licence notice: \texttt{licenses/arxiv-2307.16736-CC-BY-4.0.txt}.\par\smallskip

\noindent\textit{Editorial scope.} Counted unit: the article's lemma nonzero (source0.tex lines 1055--1057). The article's complete original proof of it is delivered with it (source0.tex lines 1058--1088, one \texttt{proof} environment of 31 lines). No mathematics is written, repaired or re-derived: the statement and the proof are the article's own lines, in the article's own notation, and no retained line is substituted or re-flowed.  No further source-local context is needed: every cross-reference of every retained line resolves inside the two delivered spans.  Nothing was omitted from any retained span, so the package carries no omission record.  No retained line contains a citation, so the wrapper carries no bibliography and no bibliography entry.  No retained line contains a figure environment, an image-inclusion command or drawing code, so no image is delivered and the package declares no asset dependency.  Editorial formatting only, in addition: an AMS \texttt{amsart} preamble that restates the article's own declarations and macros used by the retained lines, namely the environments \texttt{lemma} on separate counters, together with the standard packages those lines need.  The retained environments print in this document as Lemma 1 in order of appearance, whereas the article prints the same items with the numbering of its own full document.  The complete native source of the article as distributed in the arXiv e-print for this exact version is bundled unchanged as \texttt{sources/145010/source0.tex}.
 \begin{lemma}\label{Kloosterman sum nonzero}
 Let  $\mathfrak{N}=\tilde{s}\mathcal{O}$ be unramified with $\tilde{s}\in \mathbb{N},$  $\tilde{s}$ be squarefree. Let $\tilde{m}_1,\tilde{m}_2 \in \mathfrak{d}^{-1}_+$ and  $\omega_{\mathfrak{N}}$  be trivial. Then   $$S_{\omega_\mathfrak{N}} (\tilde{m}_1,\tilde{m}_2;1;\tilde{s})\neq 0.$$
 \end{lemma}

\begin{proof}
We have $$S_{\omega_\mathfrak{N}} (\tilde{m}_1,\tilde{m}_2;1;\tilde{s})=\prod _{\nu<\infty } S_{\omega_{\mathfrak{N},\nu}} (\tilde{m}_{1\nu},\tilde{m}_{2\nu};1;\tilde{s}_\nu).$$ Let $\tilde{s}\mathcal{O}=\prod_{l=1}^{t'}\mathfrak{p}_l$ for distinct prime ideals $\mathfrak{p}_l.$  Let $\nu_l$ be the corresponding valuation for the prime ideal $\mathfrak{p}_l.$ Hence 
$$
\prod _{\nu<\infty } S_{\omega_{\mathfrak{N},\nu}} (\tilde{m}_{1\nu},\tilde{m}_{2\nu};1;\tilde{s}_\nu)=\prod_{l=1}^{t'} S_{\omega_{\mathfrak{N},\nu_l}} (m_{1\nu_l},m_{2\nu_l};1;\varpi_{\nu_l})
$$
where $\varpi_{\nu_l}$ is a generator of the maximal ideal $(\mathfrak{p}_l)_{\nu_l}=\mathfrak{p}_l\mathcal{O}_{\nu_l}$ and   $m_{1\nu_l},m_{2\nu_l} \in  \hat{\mathfrak{d}}_{\nu_l}^{-1}.$ 
For trivial $\omega_\mathfrak{N},$ we have 
\begin{align*}
S_{\omega_{\mathfrak{N},\nu_l}} (m_{1\nu_l},m_{2\nu_l};1;\varpi_{\nu_l})&=\sum_{\substack{ s_{1},s_{2}\in \mathcal{O}_{\nu_l} / \varpi_{\nu_l} \mathcal{O}_{\nu_l}  \\{s_1s_2\equiv 1 \pmod {\varpi_{\nu_l}\mathcal{O}_{\nu_l}} }} }\theta_{\nu_l} \Big(\frac{m_{1\nu_l} s_{1}+m_{2\nu_l} s_{2}}{\varpi_{\nu_l}}\Big) \\
&=\sum_{\substack{ s_{1},s_{2}\in \mathcal{O}_{\nu_l} / \varpi_{\nu_l} \mathcal{O}_{\nu_l}  \\{s_1s_2\equiv 1 \pmod {\varpi_{\nu_l}\mathcal{O}_{\nu_l}} }} } e\Bigg(\text{Tr}\Big( \frac{m_{1\nu_l} s_{1}+m_{2\nu_l} s_{2}}{\varpi_{\nu_l}}  \Big)\Bigg),
\end{align*}
where $e(x)=e^{2\pi i x}.$
 Let $p_l=\mathfrak{p}_l \cap \mathbb{Z}$, and we have 
\begin{align*}
\Bigg[e\Bigg(\mathrm{Tr}\Big( \frac{m_{1\nu_l} s_{1}+m_{2\nu_l} s_{2}}{\varpi_{\nu_l}}  \Big)\Bigg)\Bigg]^{p_l} &= 
e\Bigg( \mathrm{Tr}\Big( p_l \cdot  \frac{m_{1\nu_l} s_{1}+m_{2\nu_l} s_{2}}{\varpi_{\nu_l}}  \Big)\Bigg) \\
&= e\Bigg( \mathrm{Tr}\Big(   \frac{m_{1  \nu_l}p_ls_{1}}{\varpi_{\nu_l}}  \Big)\Bigg)\cdot e\Bigg( \mathrm{Tr}\Big(   \frac{m_{2  \nu_l}p_ls_{2}}{\varpi_{\nu_l}}  \Big)\Bigg)\\
&= 1.
\end{align*}
The last equality follows from the fact that $ \frac{m_{1  \nu_l}p_ls_{1}}{\varpi_{\nu_l}} $ and $   \frac{m_{2  \nu_l}p_ls_{2}}{\varpi_{\nu_l}}$ belong to the local inverse different $ \hat{\mathfrak{d}}_{\nu_l}^{-1}$.

 Suppose $S_{\omega_{\mathfrak{N},\nu_l}} (m_{1\nu_l},m_{2\nu_l};1;\varpi_{\nu_l})= 0.$ 
Note that $\mathbb{Z}\Big[e^{\frac{2\pi i}{p_l}}\Big]$ is isomorphic to $\mathbb{Z}[x]/(\Phi_{p_l}(x))$ where  $e^{\frac{2\pi i}{p_l}}$ gets mapped to $x+(\Phi_{p_l}(x))$ and $\Phi_{p_l}(x)=1+x+x^2+\dots+x^{p_l-1}. $ We further have
$\mathbb{Z}[x]/(\Phi_{p_l}(x),p_l)=\mathbb{Z}[x]/((x-1)^{p_l-1},p_l)$. Consider the ring homomorphism  $$
\mathbb{Z}[x]/(\Phi_{p_l}(x))\rightarrow \mathbb{Z}[x]/((x-1)^{p_l-1},p_l) \rightarrow \mathbb{Z}[x]/((x-1),p_l)\rightarrow \mathbb{F}_{p_l},
$$
where $$x+(\Phi_{p_l}(x))\mapsto x+((x-1)^{p_l-1},p_l) \mapsto x+((x-1),p_l)$$ $$=1+[x-1]+((x-1),p_l)=1+((x-1),p_l)\mapsto 1.$$
This implies that $e^{\frac{2\pi i}{p_l}}$ gets mapped to 1 via the ring homomorphism.
Now  $e\Bigg(\mathrm{Tr}\Big( \frac{m_{1\nu_l} s_{1}+m_{2\nu_l} s_{2}}{\varpi_{\nu_l}}  \Big)\Bigg)$ lies in  $\mathbb{Z}\Big[e^{\frac{2\pi i}{p_l}}\Big]$ so that $e\Bigg(\mathrm{Tr}\Big( \frac{m_{1\nu_l} s_{1}+m_{2\nu_l} s_{2}}{\varpi_{\nu_l}}  \Big)\Bigg)$ gets mapped to $1$ via the ring homomorphism. This implies  $S_{\omega_{\mathfrak{N},\nu_l}} (m_{1\nu_l},m_{2\nu_l};1;\varpi_{\nu_l})$ will get mapped to $p_l^f-1=-1 \in \mathbb{F}_{p_l},$ where 
$\big|\mathcal{O}_{\nu_l} / \varpi_{\nu_l} \mathcal{O}_{\nu_l}\big|=p_l^f$ for some natural number $f.$ This is a contradiction since image of  $S_{\omega_{\mathfrak{N},\nu_l}} (m_{1\nu_l},m_{2\nu_l};1;\varpi_{\nu_l}) $ should be $0.$
\end{proof}

\end{document}
